\scanpage{502}
\margin{\uline{The groups} $M_{\rho,\sigma}$ \uline{etc}}

II Let us return to the situation above, forgetting (at least for the moment) $G$, $\varphi$ etc.

\margin{For a \unsure{systematic} rectification of the notations \unsure{that follow|following} cf. page 564, \uline{remark}.}

Let $\unsure{\mathcal{C}}$ be a group --- we place ourselves, for what follows, either in the category of discrete groups, or in that of profinite groups, as we choose. We shall (to fix ideas) take that of profinite groups. We suppose given
\[
(1)\qquad \rho,\sigma\in\unsure{\mathcal{C}}
\]
\margin{[\unsure{i.e. we have between} $\varepsilon_0=\ell'_0$, $\sigma=\ell'_1$ and $\rho^{-1}=\ell'_\infty$ with relation (2) $\ell'_\infty\ell'_1\ell'_0=1$]}
we put
\[
\varepsilon_0=\sigma^{-1}\rho,\qquad
\varepsilon_1=\rho\sigma^{-1}=\sigma(\varepsilon_0)=\rho(\varepsilon_0)
\]
\[
(3)\qquad
\begin{aligned}
L_0&=\text{closed subgroup generated by }\varepsilon_0\\
L_1&=\rule{6.5cm}{0.3pt}\quad\varepsilon_1
\end{aligned}
\qquad\unsure{\text{, i.}}
\]
thus $L_1=\rho(L_0)=\sigma(L_0)$.

We \deleted{suppose} consider
\diagram{\begin{tikzcd}
\mathrm{Aut}(\unsure{\mathcal{C}}) \arrow[r,"\Phi=(\underline{\unsure{\zeta|\xi}},\underline{\eta})","\deleted{\unsure{\varphi}}"'] & \deleted{\unsure{\mathcal{C}}}\,\unsure{\mathcal{C}}\times\unsure{\mathcal{C}}
\end{tikzcd}}\hfill(4)
defined by
\[
(5)\quad
\left\{\begin{aligned}
\underline{\unsure{\zeta|\xi}}(u)&=[u,\rho]=u\rho u^{-1}\rho^{-1}=u(\rho)\rho^{-1}
&&\text{thus }u(\rho)=\unsure{\zeta|\xi}\rho\quad(\unsure{\zeta|\xi}=\underline{\unsure{\zeta|\xi}}(u))\\
&=u\rho(u)^{-1}\\
\underline{\eta}(u)&=[u,\sigma]=u\sigma u^{-1}\sigma^{-1}=u(\sigma)\sigma^{-1}
&&\text{thus }u(\sigma)=\eta\sigma\quad(\eta=\underline{\eta}(u))\\
&=u\sigma(u)^{-1}.
\end{aligned}\right.
\]
Thus for $u\in\mathrm{Aut}(\unsure{\mathcal{C}})$, and $\unsure{\zeta|\xi}=\underline{\unsure{\zeta|\xi}}(u)$, $\eta=\underline{\eta}(u)$, we have
\[
\unsure{(6)|(5)}\quad
\left\{\begin{aligned}
&u\text{ fixes the }\inserted{\unsure{\mathcal{C}}}\text{-conjugation class of }\rho
\quad\Longleftrightarrow\quad \unsure{\zeta|\xi}\rho\text{ conjugate to }\rho\\
&\hspace{6cm}\Longleftrightarrow\quad\unsure{\zeta|\xi}=\alpha\rho(\alpha)^{-1}\quad(\alpha\in\unsure{\mathcal{C}})\\
&u\text{ fixes the }\unsure{\mathcal{C}}\text{-conjugation class of }\sigma
\quad\Longleftrightarrow\quad \eta\sigma\text{ conjugate to }\sigma\\
&\hspace{6cm}\Longleftrightarrow\quad\eta=\beta\sigma(\beta)^{-1}\quad(\beta\in\unsure{\mathcal{C}})\\
&u\text{ normalizes }L_0\quad\Longleftrightarrow\quad\eta^{-1}\unsure{\zeta|\xi}\in L_1
\quad\Longleftrightarrow\quad\unsure{\zeta|\xi}=\eta\varepsilon_1^\nu\quad\text{with }\nu\in\widehat{\mathbb Z}\text{ suitable}.
\end{aligned}\right.
\]
\mspage{501}

\scanpage{503}
\mspage{554}
\unsure{Let|Let} \deleted{$\unsure{B_{\rho,\sigma}}\subset$}
\[
(7)\qquad
\begin{aligned}
B_{\rho,\sigma}&\subset\mathrm{Aut}(\unsure{\mathcal{C}})\\[-2pt]
\Vert\quad&\\[-10pt]
\{u\mid{}&\text{satisfying the three conditions }(\unsure{6|5})\}
\end{aligned}
\]
\[
(8)\qquad
\begin{gathered}
\Sigma_{\rho,\sigma}\subset\unsure{\mathcal{C}}\times\unsure{\mathcal{C}}\\[-2pt]
\Vert\\[-2pt]
\left\{(\unsure{\zeta|\xi},\eta)\ \middle|\ 
\begin{aligned}
\unsure{\zeta|\xi}&=\alpha\rho(\alpha)^{-1},\quad
\eta=\beta\sigma(\beta)^{-1}\quad(\alpha,\beta\in\unsure{\mathcal{C}})\\
&\inserted{\text{i.e. }\unsure{\zeta|\xi}\rho\text{ conjugate to }\rho,\qquad\text{i.e. }\eta\sigma\text{ conjugate to }\sigma}\\
\eta^{-1}\unsure{\zeta|\xi}&\in L_1\quad\text{i.e. }\unsure{\zeta|\xi}=\eta\varepsilon_1^\nu\quad(\nu\in\widehat{\mathbb Z})
\end{aligned}\right\},
\end{gathered}
\]
so that $B_{\rho,\sigma}$ is the inverse image of $\Sigma_{\rho,\sigma}$ under (4), and we have an induced map
\diagram{\begin{tikzcd}
B_{\rho,\sigma}\arrow[r,"\varphi=\varphi_{\rho,\sigma}"]&\Sigma_{\rho,\sigma}.
\end{tikzcd}}\hfill(9)

We suppose
\[
(10)\qquad\boxed{\left\{\begin{aligned}
\text{a)}\quad&\rho,\sigma\text{ generate }\unsure{\mathcal{C}}\text{ topologically}\\
\text{b)}\quad&\text{the map (9) (injective by a))}\\
&\text{is surjective (thus bijective).}
\end{aligned}\right.}
\]
\margin{Careful, this condition \uline{is} \deleted{\unsure{not}} satisfied for $\unsure{\mathcal{C}}=\mathrm{SL}(2,\widehat{\mathbb Z})'$ ! (\unsure{it imposes} \deleted{\unsure{a condition of possibility}} $\varepsilon=\varepsilon_3(\mu)$ \unsure{on the group} $\unsure{\gamma}$ \unsure{on writing} $\unsure{\zeta|\xi},\unsure{\eta}$\dots)}

We then transport by $\varphi$ the group structure of $B_{\rho,\sigma}$ to $\Sigma_{\rho,\sigma}$. We \deleted{shall denote by}\inserted{put}
\deleted{$u_{\unsure{\zeta|\xi},\eta}=\varphi^{-1}(\unsure{\zeta|\xi},\eta)$, for}
\[
(11)\qquad u_{\unsure{\zeta|\xi},\eta}=\varphi^{-1}(\unsure{\zeta|\xi},\eta)\in B_{\rho,\sigma}
\qquad((\unsure{\zeta|\xi},\eta)\in\Sigma_{\rho,\sigma}),
\]
so that \deleted{\unsure{we shall have}} $u_{\unsure{\zeta|\xi},\eta}$ is characterized by
\[
(12)\qquad\left\{\begin{aligned}
u_{\unsure{\zeta|\xi},\eta}(\rho)&=\unsure{\zeta|\xi}\rho\\
u_{\unsure{\zeta|\xi},\eta}(\sigma)&=\eta\sigma
\end{aligned}\right.
\]
Hypothesis b) therefore means that, as soon as $(\unsure{\zeta|\xi},\eta)\in\Sigma_{\rho,\sigma}$ i.e. $\unsure{\zeta|\xi}\rho$ conjugate to $\rho$, $\eta\sigma$ conjugate to $\sigma$ and $\eta^{-1}\unsure{\zeta|\xi}\in L_1$, there exists
\mspage{502}

\scanpage{504}
an automorphism $u_{\unsure{\zeta|\xi},\eta}$ of $\unsure{\mathcal{C}}$ satisfying (12) (and which will be unique by (10) a)). The group law on $\Sigma_{\rho,\sigma}$ is made explicit as follows:
\[
(13)\quad(\unsure{\zeta|\xi}',\eta')(\unsure{\zeta|\xi},\eta)=(\unsure{\zeta|\xi}'',\eta'')\quad\text{with}\quad
\left\{\begin{aligned}
\unsure{\zeta|\xi}''&=u_{\unsure{\zeta|\xi}',\eta'}(\unsure{\zeta|\xi})\unsure{\zeta|\xi}'\\
\eta''&=u_{\unsure{\zeta|\xi}',\eta'}(\eta)\eta'.
\end{aligned}\right.
\]
Let us consider the homomorphism \deleted{$g\mapsto\mathrm{int}(g)$}
\deleted{\unsure{Let us introduce}} \deleted{Let}
\diagram{\begin{tikzcd}
\unsure{\mathcal{C}}\arrow[r]&\mathrm{Aut}(\unsure{\mathcal{C}})\\
g\arrow[r,mapsto]&\mathrm{int}(g)
\end{tikzcd}}
and the inverse image of $B_{\rho,\sigma}$ is none other than the normalizer of $L_0$ in $\unsure{\mathcal{C}}$, i.e. $N_0$. We then want to construct a group $G_{\rho,\sigma}$, containing $\unsure{\mathcal{C}}$ \inserted{as an invariant subgroup} and into which \deleted{$\unsure{B_{\rho,\sigma}}$} $B_{\rho,\sigma}$ maps, so that the tautological action of $B_{\rho,\sigma}$ on $\unsure{\mathcal{C}}$ is deduced from this homomorphism, and $G_{\rho,\sigma}$ is generated by $\unsure{\mathcal{C}}$ and the image of $B_{\rho,\sigma}$, denoted $B'_{\rho,\sigma}$, which will contain \deleted{\unsure{thus}} $N_0$.

\deleted{\unsure{To|To}}\inserted{\unsure{If}} the center of $\unsure{\mathcal{C}}$ is trivial i.e.
\diagram{\begin{tikzcd}
\unsure{\mathcal{C}}\arrow[r,hook]&\mathrm{Aut}(\unsure{\mathcal{C}})
\end{tikzcd}}
($\unsure{\mathcal{C}}$ normal subgroup),
\deleted{make this construction, let us start from
\[
B_{\rho,\sigma}\cdot\unsure{\mathcal{C}}\qquad(\text{product }\tfrac12\text{ direct})
\]
and consider the map}
\diagram{\begin{tikzcd}
\deleted{N_0}\arrow[r]&\deleted{B_{\rho,\sigma}\cdot\unsure{\mathcal{C}}}\\
\deleted{g}\arrow[r,mapsto]&\deleted{(\mathrm{int}(g),g^{-1})}
\end{tikzcd}}
it suffices to take the subgroup $G_{\rho,\sigma}$ of $\mathrm{Aut}(\unsure{\mathcal{C}})$ generated by $\unsure{\mathcal{C}}$ and $B_{\rho,\sigma}$.

\deleted{Thus $N_0$ \unsure{is|has} \unsure{commutative}, \unsure{but}}
But the hypothesis that the center of $\unsure{\mathcal{C}}$ \unsure{be|is} trivial is not a happy one --- for example it is not satisfied if $\unsure{\mathcal{C}}=\mathrm{SL}(2,\widehat{\mathbb Z})/\pm1$.
It seems to me that the introduction of $N_0$ (over which in practice we have no control)
\mspage{503}

\scanpage{505}
\mspage{555}
was not judicious either, and \deleted{that} the condition $B'_{\rho,\sigma}\supset N_0$ --- it suffices to have $B'_{\rho,\sigma}\supset L_0$.
I shall therefore consider
\[
(14)\quad\left\{\begin{gathered}
S_0G_{\rho,\sigma}\overset{\mathrm{dfn}}{=}B_{\rho,\sigma}\cdot\unsure{\mathcal{C}}\qquad\text{the product }\tfrac12\text{ direct, and}
\end{gathered}\right.
\]
\diagram{\begin{tikzcd}
L_0\arrow[r,hook,"i"]&\deleted{B_{\rho,\sigma}\cdot\unsure{\mathcal{C}}}\ S_0G_{\rho,\sigma}\\
g\arrow[r,mapsto]&\mathrm{int}(g^{\unsure{\mathcal{C}}})\cdot g^{-1}
\end{tikzcd}}
the image of this homomorphism is a \unsure{normal} subgroup, for a) it commutes with $\unsure{\mathcal{C}}$ (trivial) and b) it is invariant under $B_{\rho,\sigma}$, for we shall have for $u\in B_{\rho,\sigma}$
\[
\mathrm{int}(u)(i(g))=\mathrm{int}(u)(\mathrm{int}(g)\cdot g^{-1})
=\mathrm{int}(u(g))\cdot\underbrace{\mathrm{int}(u)(g^{-1})}_{u(g^{-1})=u(g)^{-1}}
=i(u(g))
\]
(i.e. $i$ commutes with the actions of $B_{\rho,\sigma}$, operating on $L_0$ by restriction to $L_0$ of the \inserted{\unsure{tautological|total}} operation of $B_{\rho,\sigma}$ on $\unsure{\mathcal{C}}$, and on \deleted{$B_{\rho,\sigma}\cdot\unsure{\mathcal{C}}$}\inserted{$S_0G_{\rho,\sigma}$} $\supset B_{\rho,\sigma}$ by inner automorphisms).

We can therefore consider
\[
(15)\qquad G_{\rho,\sigma}\overset{\mathrm{dfn}}{=}S_0G_{\rho,\sigma}/i(L_0)
\]
then the inclusion homomorphisms
\diagram{\begin{tikzcd}
B_{\rho,\sigma}\arrow[r,hook]&S_0G_{\rho,\sigma}
\end{tikzcd}},
\diagram{\begin{tikzcd}
\unsure{\mathcal{C}}\arrow[r]&S_0G_{\rho,\sigma}
\end{tikzcd}}
give, by composition with
\diagram{\begin{tikzcd}
S_0G_{\rho,\sigma}\arrow[r]&G_{\rho,\sigma}
\end{tikzcd}},
homomorphisms
\diagram{\begin{tikzcd}
\deleted{\unsure{B_{\rho,\sigma}}}\unsure{\mathcal{C}}\arrow[r,"{\deleted{\unsure{\lhook}}}"{pos=0.08}]&G_{\rho,\sigma}
\end{tikzcd}},
\diagram{\begin{tikzcd}
B_{\rho,\sigma}\arrow[r,hook]&G_{\rho,\sigma}
\end{tikzcd}}\hfill(16)
making the diagram commutative
\diagram{\begin{tikzcd}[column sep=small,row sep=small]
&\unsure{\mathcal{C}}\arrow[dr,"{\deleted{\unsure{\lhook}}}"{pos=0.08}]&\\
L_0\arrow[ur,hook]\arrow[dr]&&G_{\rho,\sigma}\\
&B_{\rho,\sigma}\arrow[ur,hook]&
\end{tikzcd}}\hfill(17)
We note
\[
(18)\qquad\mathrm{Ker}(\unsure{\mathcal{C}}\longrightarrow G_{\rho,\sigma})
=L_0\cap\mathrm{Centr}(\unsure{\mathcal{C}})
=\mathrm{Ker}(L_0\longrightarrow B_{\rho,\sigma})
\]
\mspage{504}

\scanpage{506}
In the case $\unsure{\mathcal{C}}=\mathrm{SL}(2,\widehat{\mathbb Z})'$, (with $\rho,\sigma$ as usual) we shall have $L_0\cap\mathrm{Centr}(\unsure{\mathcal{C}})=\{1\}$. Thus it does not seem a priori that there are any major drawbacks to supposing that
\[
(18)\qquad L_0\cap\mathrm{Centr}(\unsure{\mathcal{C}})=\{1\}\quad\text{i.e.}
\]
\diagram{\begin{tikzcd}
\unsure{\mathcal{C}}\arrow[r,hook]&\unsure{G}_{\rho,\sigma}\quad\text{injective}
\end{tikzcd}}
Let us note that \deleted{$L_0$ \unsure{is going to be|is its}} the image of
\diagram{\begin{tikzcd}[row sep=small]
L_0\arrow[r]&B_{\rho,\sigma}\\
\ell\arrow[r,mapsto]&\operatorname{int}(\ell)
\end{tikzcd}}
(composite of (14) and the projection $S_0G_{\rho,\sigma}\longrightarrow B_{\rho,\sigma}$) is obviously \unsure{invariant|invariant}, and putting
\[
(19)\qquad\unsure{\Gamma_{\rho,\sigma}^{*}|\gamma_{\rho,\sigma}}=B_{\rho,\sigma}/\operatorname{Im}L_0
\]
we find
\[
(20)\qquad G_{\rho,\sigma}/\unsure{\mathcal{C}}_{\rho,\sigma}\simeq B_{\rho,\sigma}/L_0\simeq\unsure{\Gamma|\gamma}_{\rho,\sigma}.
\]
\margin{i.e. homomorphism of exact sequences
\diagram{\begin{tikzcd}[column sep=small,row sep=small]
1\arrow[r]&\deleted{\unsure{\mathcal{C}_{\rho,\sigma}}}\,L_0\arrow[r]\arrow[d]&B_{\rho,\sigma}\arrow[r]\arrow[d]&\unsure{\Gamma|\gamma}_{\rho,\sigma}\arrow[r,dash]\arrow[d,equal]&1\\
1\arrow[r]&\unsure{\mathcal{C}}\arrow[r]\arrow[d,equal]&G_{\rho,\sigma}\arrow[r]&\unsure{\Gamma|\gamma}_{\rho,\sigma}\arrow[r,dash]&1\\
&\unsure{SG_{\rho,\sigma}}&&&
\end{tikzcd}}(20bis)}

(\uline{N.B.} In the case $\unsure{\mathcal{C}}=\mathrm{SL}(2,\widehat{\mathbb Z})'$ ; \deleted{$\unsure{B_{\rho,\sigma}}$} \unsure{etc.|and}, $B_{\rho,\sigma}$ \inserted{$\subset\unsure{\mathrm{GP}}(\widehat{\mathbb Z})$} is the group previously denoted $B'_0$ (or rather its isomorphic image $\widetilde B'_0$), and the homomorphism $\mathrm{dét}\mid B'_0=B_{\rho,\sigma}$ \deleted{induces} (trivial on $L_0$) induces by passage to the quotient an \uline{isomorphism}
\diagram{\begin{tikzcd}
\unsure{\Gamma|\gamma}_{\rho,\sigma}\arrow[r,"\sim"]&G_{\mathrm m}
\end{tikzcd}}\hfill(21)
\margin{In the case where we start from $\mathrm{SL}(2,\widehat{\mathbb Z}),\rho,\sigma$, we end up with $G=\mathrm{GL}(2,\widehat{\mathbb Z})$\dots}
and we see, \deleted{by} using the isomorphism
\[
B'_0=\unsure{\widetilde T_0|\widetilde\Gamma_0}\cdot L_0\quad(\tfrac12\text{ direct})
\]
\inserted{where $\unsure{\widetilde T_0|\overline T_0}\xrightarrow{\sim}B'_0/L_0\simeq G_{\mathrm m}$}
\[
\text{i.e.}\qquad B_{\rho,\sigma}\simeq\unsure{\Gamma|\gamma}_{\rho,\sigma}\cdot L_0,
\]
that in this case we have
\mspage{\unsure{505}}
\[
(22)\qquad G_{\rho,\sigma}=\deleted{\unsure{\mathrm{GL}}}\mathrm{GL}(2,\widehat{\mathbb Z})'.
\]
\mspage{505}